% =============================================================================
\section{How the results were reached}
\label{app:provenance}

The work was kept in a running findings ledger and a task log, with corrections recorded as
they happened. It began on 12~October~2026, when I asked whether the tetrahedron has any
place in the theory of Jordan algebras. By 17~October it had become a study of quantum
tetrahedra over the four rank-three Jordan algebras, in which the volume of a quantum
tetrahedron reappears as the commutator of two pair Casimirs. From there it moved to the
dynamics of loop quantum gravity: whether a Lorentzian Hamiltonian constraint can close on
an off-shell space. That question was not settled, and the companion
paper~\cite{FitchVolumeGap} describes it. This appendix follows the branch that became the
present paper.

On 20~October, testing whether self-dual variables could tell Lorentzian dynamics apart from
a formal factor of $i$, the four-valent volume with $\mathfrak{su}(1,1)$ labels was written
as a tridiagonal matrix and found not to be essentially self-adjoint. The first reading was
that the extension angle was an invisible phase. A referee corrected it: for equal labels on
two legs only $t=0$ and $t=\frac\pi2$ respect the orientation, and they give different
smallest volumes. A literature search showed that the closed form of the matrix elements was
classical, and that the infinite representations behind it had been studied. What was not
found anywhere was the limit-circle property and the extension it forces, so that was proved
first, for four classes of chains. The novelty assessment was candid: probably new, but a
short exercise for a specialist.

The same afternoon the spectra of these extensions were seen to obey a Bohr--Sommerfeld rule
in which the extension angle appears as a phase. A referee found the first account's claim
about the meaning of $t=0$ and $t=\frac\pi2$ wrong, since it depends on how the legs are
labelled. The rigorous proof rests on the factorisation of the transfer matrix into shears,
which is classical~\cite{Kac1999}, and on an exact rotation identity for two-site blocks.
With these it reached the counting law for cubic links, but at first without a rate. The
rate $O(\lambda^{-1/3})$ was first derived from Geronimo's turning-point
theorem~\cite{Geronimo2009}. A referee then found that the printed proof of that theorem
relies on an Airy weight that cannot exist, so the rate was only conditional. It was
replaced by the discrete Olver bound of \Cref{thm:olver}, proved from scratch, and a second
referee judged the rate unconditional. Meanwhile the rule was proved for every exponent
$\alpha>1$, and the second-order constant $(2\alpha-1)/24$, first seen numerically, was
proved as well. A check of this stage by Codex confirmed the cubic rate. It also found an
auxiliary threshold that was false, with the counterexample $a_n=n^2+100n$.

A novelty audit of the general law found it new but not large: its leading constant is the
one Valent's conjecture predicts, and the Codex check located isolated leading terms for
special cubic chains in earlier work cited by Gilewicz, Leopold, Ruffing and
Valent~\cite{GLRV2006}. I judged that a general law with a predicted constant was not
enough, and asked for named open problems. That led to three of them. The first was the
dimerised links of \Cref{sec:dimer} and Valent's conjecture, where a referee showed that the
indicator follows in half a page from known results, so that what is new is the $O(1)$
term. The second was the question of Fischer, Burgarth and Lonigro~\cite{FBL2026}, answered
in \Cref{thm:fbl}. The third was the open problem of Berg and Szwarc~\cite{BergSzwarc2026}
on critical indices. Our determinacy criterion for it turned out to be, in substance, a
theorem of Borichev and Sodin~\cite{BorichevSodin1998}, and the attribution was checked
against their 1998 preprint before it was cited. That material was first planned as a
separate paper and was folded into \Cref{sec:moments} on 21~October, because it does not
stand on its own.

The paper was written from 21~October, from the primary sources and not from our working
notes. I asked for a revision that explains the central mechanisms, the phase budget behind
the missing quarter-period and the Bessel far field, before their proofs, with the routine
estimates moved to \Cref{app:routine}. An independent Claude referee then read the whole
paper and checked its constants with its own Wronskian code. It found no wrong theorem, but
it did find a gap in a remark on the class~$\mathbf A$, which was repaired. A GPT referee
found that the reflection formula for negative eigenvalues in \Cref{cor:squeeze} was off by
one in its index. It also found that the solutions defining the boundary conditions must
solve the recurrence only away from the first site. Both findings were confirmed against
\Cref{thm:count} before they were adopted. The final revision also found and fixed an
inconsistent diagonal value on one page. When the code was first published, 52 of its
checksums failed because the repository had normalised line endings; this was fixed before
the release was cited here.
